\documentclass[11pt]{amsart}

\usepackage[utf8]{inputenc}
\usepackage{amsmath,amssymb,amsthm,mathtools}
\usepackage[margin=1.05in]{geometry}
\usepackage{enumitem}
\usepackage{microtype}
\usepackage{booktabs}
\usepackage{xcolor}
\usepackage{listings}
\usepackage[colorlinks=true,linkcolor=blue!50!black,citecolor=blue!50!black,urlcolor=blue!50!black]{hyperref}

\lstset{basicstyle=\ttfamily\scriptsize,breaklines=true,columns=fullflexible,
  frame=single,framerule=0.3pt,rulecolor=\color{black!40},showstringspaces=false,
  keepspaces=true,upquote=true}

\newtheorem{theorem}{Theorem}
\renewcommand{\thetheorem}{\Alph{theorem}}
\newtheorem{maincor}[theorem]{Corollary}
\newtheorem{lemma}{Lemma}[section]
\newtheorem{proposition}[lemma]{Proposition}
\newtheorem{corollary}[lemma]{Corollary}
\newtheorem{thm}[lemma]{Theorem}
\theoremstyle{definition}
\newtheorem{definition}[lemma]{Definition}
\newtheorem{problem}[lemma]{Problem}
\newtheorem{principle}[lemma]{Principle}
\newtheorem{question}[lemma]{Question}
\theoremstyle{remark}
\newtheorem{remark}[lemma]{Remark}
\newtheorem{example}[lemma]{Example}

\newcommand{\N}{\mathbb{N}}
\newcommand{\Z}{\mathbb{Z}}
\newcommand{\Q}{\mathbb{Q}}
\newcommand{\R}{\mathbb{R}}
\newcommand{\C}{\mathbb{C}}
\newcommand{\fl}[1]{\lfloor #1\rfloor}
\newcommand{\fr}[1]{\{#1\}}

\title[Carry forcing and the Erd\H{o}s--Graham prediction]{Carry Forcing for Zero--One
Relations:\\ the Erd\H{o}s--Graham Completeness Prediction\\ for the Bases
$\gamma^k=\gamma^a+\gamma^b$ below $\sqrt2$}

\author{Keio}
\thanks{Research draft, October 29, 2026. The results and the exposition were developed
interactively with an AI assistant (Claude, Anthropic); the author is responsible for the
content. All numerical values quoted in the paper are reproduced by the script in
Appendix~\ref{app:script}. This draft has not been peer reviewed.}

\date{October 29, 2026}

\begin{document}

\begin{abstract}
For $t>0$ and $\gamma>1$ let $s_n=\fl{t\gamma^n}$ $(n\ge 0)$. A sequence of non-negative
integers is \emph{complete} if every sufficiently large integer is a sum of distinct terms.
Graham asked for which pairs $(t,\gamma)$ the sequence $(s_n)$ is complete, and Erd\H{o}s and
Graham wrote that it is thought to be complete for every $t>0$ whenever
$1<\gamma<\varphi=\frac{1+\sqrt5}{2}$ (Erd\H{o}s Problem \#349). Geneson recently refuted this
for a Salem base; on the positive side, completeness for all $t>0$ at the base
$\gamma=\sqrt\varphi$ follows from a recent Lean-verified argument of Kitamura.

We formulate a \emph{carry-forcing problem}: if $\gamma^k=\sum_{i\in I}\gamma^i$ with
$I\subseteq\{0,\dots,k-1\}$, then $s_{n+k}=\sum_{i\in I}s_{n+i}+\delta_n$ with a ``carry''
$\delta_n\in\{0,\dots,|I|-1\}$; is $\delta_n\neq0$ infinitely often for every $t>0$? Combining
the irrationality of $\gamma$, Descartes' rule of signs and a sign-change lemma for
exponential polynomials, we prove that the answer is always yes (for the golden ratio this is
Kitamura's observation, and for Pisot trinomial bases it can also be derived from a theorem
of Dubickas). For trinomial relations $\gamma^k=\gamma^a+\gamma^b$ the carries are units, so
every floor orbit contains infinitely many disjoint \emph{unit switches}
$\{s_{n+a},s_{n+b}\}\leftrightarrow\{s_{n+k}\}$. Our main new tool is a switch-extraction lemma,
valid exactly below $\sqrt2$, which removes these switches from a single orbit without
destroying Graham's interval propagation. We obtain:
\begin{enumerate}[label=(\roman*),leftmargin=2em]
\item if $\gamma\in(1,\sqrt2)$ and $\gamma^k=\gamma^a+\gamma^b$ with $k>a>b\ge0$ --- an infinite
set of bases, accumulating at $1$ and containing the two smallest Pisot numbers --- then
$(\fl{t\gamma^n})_{n\ge0}$ is strongly complete for \emph{every} $t>0$. When the reduced
trinomial $x^m-x^j-1$ $(m=k-b,\ j=a-b)$ has $\gcd(m,j)\ge2$, this also follows from (ii) by
splitting the orbit (Kitamura's device); the extraction lemma handles the primitive case
$\gcd(m,j)=1$, which includes the plastic number and the root of $x^4-x^3-1$;
\item for every $\gamma\in(1,2)$ with $\gamma^k=\gamma^a+\gamma^b$, including $\varphi$ and
$1.4655\ldots$, and all $\alpha,\beta>0$, the union of $(\fl{\alpha\gamma^n})$ and
$(\fl{\beta\gamma^n})$ is complete. For $\gamma=\varphi$ this was observed before, and (ii)
answers the existential form of Erd\H{o}s Problem \#354(ii) recorded in the Formal Conjectures
database, which Kitamura had already answered.
\end{enumerate}
The threshold $\sqrt2$ is sharp for the extraction step. We also explain why base $2$ admits
no exact relations of this kind.
\end{abstract}

\maketitle

\tableofcontents

%======================================================================
\section{Introduction}\label{sec:intro}
%======================================================================

\subsection{Complete sequences of integer parts of powers}
A sequence $(q_i)_{i\in\mathcal I}$ of non-negative integers (indexed, so that equal
values carried by different indices may both be used) is \emph{complete} if every
sufficiently large integer is of the form $\sum_{i\in J}q_i$ with $J\subseteq\mathcal I$
finite; it is \emph{strongly complete} if it stays complete after deleting any finite set
of terms. For $t>0$ and $\gamma>1$ we write
\[
  S_t(\gamma)=(s_n)_{n\ge 0},\qquad s_n=\fl{t\gamma^n},\qquad f_n=\fr{t\gamma^n}\in[0,1).
\]
Graham asked for which pairs $(t,\gamma)$ the sequence $S_t(\gamma)$ is complete; according to
\cite{vD26} he determined the answer for $0<t<1$, $1<\gamma<2$, and even there the set of good
pairs is very complicated \cite{EP349}. Erd\H{o}s and Graham \cite[p.~57]{EG80} wrote that
$S_t(\gamma)$ is thought to be complete for all $t>0$ when $1<\gamma<\varphi=\frac{1+\sqrt5}2$;
in the formulation of Erd\H{o}s Problem \#349 \cite{EP349}, ``it seems likely that the
sequence is complete for all $t>0$ and all $1<\alpha<\frac{1+\sqrt5}{2}$''. Van Doorn
\cite{vD26} settled the range $\gamma\ge\varphi$ (for instance, for
$\varphi\le\gamma\le5^{1/3}$ the sequence is complete if and only if
$t<\min(3/\gamma^2,5/\gamma^3)$), and for $1<\gamma<\varphi$ he proved completeness in several
bounded ranges of $t$ (partly computer-assisted, e.g.\ $1.3<\gamma\le1.4$, $t\le3$) and in the
region $1<\gamma\le1+1/(\lceil t\rceil+2\lceil\sqrt t\rceil)$; large $t$ remained open there.
Geneson \cite[Theorem~9]{Gen26} then showed that the prediction fails in general: there are a
Salem number $\gamma\in(\frac65,\frac{13}{10})$ and arbitrarily large $t$ for which every term of
$S_t(\gamma)$ is even.

Erd\H{o}s Problem \#354 \cite{EP354}, first mentioned by Graham in 1971 (see \cite{EP354}),
asks: \emph{if $\alpha,\beta>0$ and $\alpha/\beta$ is irrational, is the multiset
$\{\fl{2^n\alpha}\}_{n\ge0}\cup\{\fl{2^n\beta}\}_{n\ge0}$ complete? What if $2$ is replaced
by some $\gamma\in(1,2)$?} The Formal Conjectures database \cite{FC} records part (ii) as the
existential statement
\begin{equation}\label{eq:354ii}
  \exists\,\gamma\in(1,2)\ \ \forall\,\alpha,\beta>0\ \text{with}\ \alpha/\beta\notin\Q:\quad
  \{\fl{\gamma^n\alpha}\}_{n\ge0}\cup\{\fl{\gamma^n\beta}\}_{n\ge0}\ \text{is complete}
\end{equation}
(see Section~\ref{sec:proofC} for the exact formal statement). Geneson
\cite[Corollary~12]{Gen26} constructed, for his Salem base, two such sequences with irrational
ratio whose union is incomplete, so \eqref{eq:354ii} cannot hold for \emph{every}
$\gamma\in(1,2)$.

\subsection{Prior work that this paper builds on}\label{sec:prior}
Kitamura \cite{Kit26} gave a Lean-verified affirmative answer to \eqref{eq:354ii}, with
$\gamma=\sqrt\varphi$, in a public repository dated September 5, 2026 (proposed to \cite{FC} in
pull request \#5286, not merged as of October 28, 2026). His argument is the golden-ratio case of
the mechanism used here: for $c_n=\fl{u\varphi^n}$ one has $c_{n+2}=c_n+c_{n+1}+\delta_n$ with
$\delta_n\in\{0,1\}$; $\delta_n=1$ occurs arbitrarily far out, because otherwise the fractional
parts would be a bounded non-negative solution of the Fibonacci recurrence and hence zero; the
resulting disjoint ``carry-one triples'' produce long runs of consecutive subset sums, which are
propagated by the odd-indexed terms $\fl{(u\sqrt\varphi)\varphi^n}$. In particular his argument
shows that $S_t(\sqrt\varphi)$ is complete for every $t>0$. His repository also mentions an earlier
proof note for the root of $\gamma^{10}=\gamma+1$ \cite{Kit26}, which we have not been able to
locate; we do not know its argument. An automated review of the formal statement (record of
September 23, 2026 in \cite{Rev26}, and formal-conjectures issue \#6542) gave the same argument
with $\gamma=\varphi$ and arbitrary $\alpha,\beta>0$ --- this is Theorem~\ref{thm:C} below for
$\gamma=\varphi$ --- and argued that the existential statement \eqref{eq:354ii} does not capture
Graham's question, which asks about a general base $\gamma$. For Pisot bases, the conclusion of
Theorem~\ref{thm:carrygeneral}(iii) below can be deduced from a theorem of Dubickas \cite{Dub06}
(Remark~\ref{rem:dub}). Another review record in \cite{Rev26} attributes to Hegyv\'ari
\cite[Theorem~4]{He89} a completeness theorem for unions of three base-$2$ sequences; if it
applies as quoted there, completeness of $\fl{t\,2^{n/k}}$ for all $t>0$ and $k\ge3$ would follow.
We have not been able to verify this. Finally, Yu and Chen \cite{YC26} have announced a
Lean-verified affirmative answer to part (i) of Problem \#354 (base $2$), proposed to \cite{FC}
in pull request \#6433. We have not checked these works and make no claim about them beyond what
is stated here.

\subsection{The carry-forcing problem}
Every exact additive relation among powers of $\gamma$ survives in the floor orbit up to a
bounded error. If
\begin{equation}\label{eq:rel}
  \gamma^k=\sum_{i\in I}\gamma^i,\qquad I\subseteq\{0,1,\dots,k-1\},\ |I|\ge2,
\end{equation}
then $t\gamma^{n+k}=\sum_{i\in I}t\gamma^{n+i}$, and taking integer parts,
\begin{equation}\label{eq:carry}
  s_{n+k}=\sum_{i\in I}s_{n+i}+\delta_n,\qquad
  \delta_n=\Bigl\lfloor\sum_{i\in I}f_{n+i}\Bigr\rfloor\in\{0,1,\dots,|I|-1\}.
\end{equation}
We call $\delta_n=\delta_n(t)$ the \emph{carry} at $n$.

\begin{problem}[Carry forcing]\label{prob:carry}
Let $\gamma>1$ satisfy a relation \eqref{eq:rel}. Is it true that for \emph{every} $t>0$
the carry $\delta_n(t)$ is non-zero for infinitely many $n$? When are the carries units?
\end{problem}

When $\delta_n=1$ the two disjoint index sets $\{n+i:i\in I\}$ and $\{n+k\}$ carry sums that
differ by exactly $1$, a \emph{unit switch}. Many disjoint unit switches produce long runs of
consecutive subset sums, which can be propagated to all large integers by Graham's interval
argument, provided that the terms reserved for the switches are not needed by the propagation.
For two sequences (or for $\gamma=\sqrt\varphi$, whose orbit splits into two golden-ratio orbits)
the second sequence propagates; for a \emph{single} orbit the switches must be extracted from the
propagating sequence itself, and this is where the threshold $\sqrt2$ appears.

\subsection{Main results}
We call $\gamma>1$ a \emph{trinomial base} if $\gamma^k=\gamma^a+\gamma^b$ for some integers
$k>a>b\ge0$, i.e.\ if $\gamma$ is the positive root of $x^k-x^a-x^b$; ``trinomial'' always refers
to polynomials of this shape (other trinomials, such as $x^4-2x-1$, are not covered).

\begin{theorem}[Carry forcing]\label{thm:A}
Let $k>a>b\ge0$ be integers and let $\gamma>1$ be a root of $x^k-x^a-x^b$. Then
$1<\gamma<2$, $\gamma$ is irrational, and for every $t>0$
\[
  \delta_n:=\fl{t\gamma^{n+k}}-\fl{t\gamma^{n+a}}-\fl{t\gamma^{n+b}}\in\{0,1\}
  \quad\text{for all }n\ge0,
\]
and $\delta_n=1$ for infinitely many $n$.
More generally, for every relation \eqref{eq:rel} and every $t>0$ the carries
\eqref{eq:carry} are non-zero for infinitely many $n$ (Theorem~\ref{thm:carrygeneral}).
\end{theorem}

\begin{theorem}[The Erd\H{o}s--Graham prediction for trinomial bases below $\sqrt2$]\label{thm:B}
Let $\gamma\in(1,\sqrt2)$ be a root of a trinomial $x^k-x^a-x^b$ with $k>a>b\ge0$. Then for
every $t>0$ the sequence $(\fl{t\gamma^n})_{n\ge0}$ is complete; in fact it is strongly
complete, and the set of its values is complete.
\end{theorem}

By Lemma~\ref{lem:trinroots}, the bases covered by Theorem~\ref{thm:B} are exactly the
positive roots of $x^m-x^j-1$ $(m>j\ge1)$ other than $\varphi$ (root of $x^2-x-1$) and
$1.46557\ldots$ (root of $x^3-x^2-1$). They form an infinite set accumulating only at $1$,
whose two largest elements are $1.38027\ldots$ (root of $x^4-x^3-1$) and the plastic number
$1.32471\ldots$ (root of $x^3-x-1$), the two smallest Pisot numbers \cite{Sie44}; see
Table~\ref{tab:roots}. The base $\sqrt\varphi=1.27201\ldots$ (root of $x^4-x^2-1$) of
\cite{Kit26} and the root $1.07576\ldots$ of $x^{10}-x-1$ are among them.

\begin{theorem}[Unions for all trinomial bases]\label{thm:C}
Let $\gamma\in(1,2)$ be a root of a trinomial $x^k-x^a-x^b$ with $k>a>b\ge0$ (for instance
$\gamma=\varphi$). Then for all $\alpha,\beta>0$ the union of
$(\fl{\alpha\gamma^n})_{n\ge0}$ and $(\fl{\beta\gamma^n})_{n\ge0}$ is complete, indeed
strongly complete.
\end{theorem}

Theorem~\ref{thm:C} is a direct extension of the argument of \cite{Kit26,Rev26} from the golden
ratio to all trinomial bases. It gives the following consequences.

\begin{maincor}[Problem \#354(ii), existential form]\label{cor:D}
Statement \eqref{eq:354ii}, in the exact form recorded in \cite{FC}, holds with $\gamma$ equal to
any root in $(1,2)$ of a trinomial $x^k-x^a-x^b$ (Kitamura \cite{Kit26}: $\gamma=\sqrt\varphi$).
The conclusion holds for all $\alpha,\beta>0$, without the hypothesis $\alpha/\beta\notin\Q$, and
for the trinomial roots $\gamma<\sqrt2$ the union is complete even when regarded as a set of
values rather than as a multiset.
\end{maincor}

\begin{maincor}[Multiplicity two at the golden ratio]\label{cor:E}
For every $t>0$, every sufficiently large integer $N$ can be written as
$N=\sum_{n\ge0}c_n\fl{t\varphi^n}$ with $c_n\in\{0,1,2\}$. By \cite{vD26}, for
$t\ge 3/\varphi^2=1.1459\ldots$ the coefficients cannot in general be restricted to
$\{0,1\}$.
\end{maincor}

\subsection{What is new}
To our knowledge the following are new: the general carry-forcing theorem for $0/1$-relations
(Theorem~\ref{thm:carrygeneral}), whose proof replaces the explicit Binet formula of the
golden-ratio case by Descartes' rule of signs and a sign-change lemma for exponential
polynomials (for Pisot roots its conclusion can also be derived from \cite{Dub06}, see
Remark~\ref{rem:dub}); the switch-extraction lemma (Lemma~\ref{lem:extract}) together with the sharpness
of its threshold $\sqrt2$ (Proposition~\ref{prop:sharp}); and, as a consequence,
Theorem~\ref{thm:B}, which confirms the Erd\H{o}s--Graham prediction for the infinite family of
all trinomial bases below $\sqrt2$. For the trinomials $x^m-x^j-1$ with $\gcd(m,j)\ge2$ (for
instance $x^4-x^2-1$, whose root is $\sqrt\varphi$) Theorem~\ref{thm:B} also follows from
Theorem~\ref{thm:C} by splitting the orbit into $\gcd(m,j)$ orbits of a trinomial base
(Remark~\ref{rem:split}); this is Kitamura's device. The switch-extraction lemma is used for the
primitive trinomials, $\gcd(m,j)=1$, an infinite subfamily which contains the two smallest Pisot
numbers. Theorem~\ref{thm:C} and Corollaries~\ref{cor:D} and~\ref{cor:E} are routine extensions
of the argument of \cite{Kit26,Rev26} (for $\gamma=\varphi$, Theorem~\ref{thm:C} is stated in
\cite{Rev26}) and are included for completeness.

\subsection{The principle behind the proofs}
\begin{principle}[Carry-forced switches]\label{prin}
An exact $0/1$-relation among powers of an irrational algebraic integer $\gamma$ survives
in every floor orbit $\fl{t\gamma^n}$ up to a carry. If the carries vanished from some
point on, the fractional parts $\fr{t\gamma^n}$ would form a bounded non-negative solution of
the linear recurrence attached to the relation. Because the relation polynomial has $\gamma$
as its only positive root (Descartes), boundedness removes the $\gamma$-component, and a
non-zero real exponential polynomial without positive real characteristic roots takes
negative values; so the fractional parts would vanish, forcing $\gamma\in\Q$. Hence carries
occur infinitely often. For trinomials they are unit switches, and below $\sqrt2$ they can
be removed from the orbit without destroying the growth condition that drives Graham's
interval propagation.
\end{principle}

\subsection{What is not claimed}
\begin{itemize}[leftmargin=1.5em]
\item We do not claim to be the first to answer \eqref{eq:354ii}; see Section~\ref{sec:prior},
where it is also noted that \eqref{eq:354ii} has been criticised as a misformalisation of
Graham's question. Nothing is claimed about part (i) of Problem \#354 (base $2$). Distinct sets of powers of $2$
have distinct sums, so the mechanism of this paper is unavailable there
(Section~\ref{sec:disc}).
\item The characterisation asked for in Problem \#349 remains open. Theorem~\ref{thm:B}
settles the vertical lines $\{\gamma\}\times(0,\infty)$ of the $(\gamma,t)$-plane for the
trinomial roots $\gamma<\sqrt2$. For the trinomial roots $\varphi$ and $1.46557\ldots$ the
single-sequence question is not settled by our method (for $\varphi$ completeness fails for
$t\ge3/\varphi^2$ by \cite{vD26}).
\item The threshold $\sqrt2$ is sharp for the extraction step only; it is not known to be sharp
for completeness.
\end{itemize}

\subsection*{Organisation}
Section~\ref{sec:prelim} contains the propagation lemma, a growth bound and a sign-change
lemma for exponential polynomials. Section~\ref{sec:carry} proves carry forcing
(Theorem~\ref{thm:A} and its generalisation) and describes the trinomial bases.
Section~\ref{sec:extract} proves the switch-extraction lemma and its sharpness.
Theorems~\ref{thm:B} and~\ref{thm:C} and the corollaries are proved in
Sections~\ref{sec:proofB} and~\ref{sec:proofC}. Section~\ref{sec:disc} discusses base $2$,
Salem bases and open problems, and Section~\ref{sec:num} summarises the numerical
verification.


%======================================================================
\section{Preliminaries}\label{sec:prelim}
%======================================================================

\begin{lemma}[Interval propagation]\label{lem:prop}
Let $(q_m)_{m\ge0}$ be non-negative integers with $\sum_m q_m=\infty$, and let $L\ge1$ be an
integer with
\begin{equation}\label{eq:propcond}
  q_m\le L+\sum_{i<m}q_i\qquad(m\ge0).
\end{equation}
Let $\mathcal P$ be a finite family of non-negative integers, disjoint (as an indexed
family) from $(q_m)$, whose subset sums contain the $L$ consecutive integers
$u,u+1,\dots,u+L-1$. Then every integer $N\ge u$ is a sum of distinct members of
$\mathcal P\sqcup(q_m)_{m\ge0}$.
\end{lemma}

\begin{proof}
Put $I_m=[u,\,u+L-1+\sum_{i<m}q_i]\cap\Z$. Every element of $I_0$ is a subset sum of
$\mathcal P$. Suppose every element of $I_m$ is a sum of distinct members of
$\mathcal P\sqcup\{q_0,\dots,q_{m-1}\}$. Then so is every element of $I_m+q_m$ (add the
new term $q_m$), and $I_m\cup(I_m+q_m)=I_{m+1}$ because $q_m\le|I_m|=L+\sum_{i<m}q_i$ by
\eqref{eq:propcond}. Since $\sum q_m=\infty$, $\bigcup_mI_m=[u,\infty)\cap\Z$.
\end{proof}

This is the mechanism behind Brown's criterion \cite{Bro61} and behind Graham's theorem that
the union of a ``$\Sigma$-sequence'' (one with $q_m-\sum_{i<m}q_i$ bounded above) and a
``precomplete'' sequence (one whose subset sums contain arbitrarily long runs of
consecutive integers) is complete \cite{Gra64}.

\begin{lemma}[Growth below $2$]\label{lem:growth}
Let $1<\gamma<2$, $t>0$, $s_n=\fl{t\gamma^n}$, and put
\[
  \rho=\rho(\gamma)=\frac{2-\gamma}{\gamma-1}>0,\qquad D(m)=s_m-\sum_{i<m}s_i .
\]
Then $D(m)\le -t\rho\gamma^m+\frac{t}{\gamma-1}+m$ for all $m\ge0$. In particular
$c_0(t,\gamma):=\max_{m\ge0}D(m)$ exists, i.e.\ $S_t(\gamma)$ is a $\Sigma$-sequence.
\end{lemma}

\begin{proof}
$s_m\le t\gamma^m$ and $s_i\ge t\gamma^i-1$, so
$D(m)\le t\gamma^m-t\frac{\gamma^m-1}{\gamma-1}+m=-t\rho\gamma^m+\frac t{\gamma-1}+m$.
The right-hand side tends to $-\infty$, and $D(m)\in\Z$.
\end{proof}

We use the following standard facts on exponential polynomials; we include a proof for
completeness (compare \cite{BG07} for finer statements on the sign patterns of linear
recurrence sequences).

\begin{lemma}[Exponential polynomials]\label{lem:exppoly}
Let $u_n=\sum_{j=1}^rP_j(n)\lambda_j^n$ $(n\ge n_0)$ be real numbers, where $r\ge1$,
$\lambda_1,\dots,\lambda_r\in\C\setminus\{0\}$ are distinct and $P_1,\dots,P_r\in\C[x]$ are
non-zero polynomials.
\begin{enumerate}[label=(\alph*)]
\item If $(u_n)$ is bounded, then $|\lambda_j|\le1$ for all $j$, and $\deg P_j=0$ whenever
$|\lambda_j|=1$.
\item If no $\lambda_j$ is a positive real number, then $u_n<0$ for infinitely many $n$.
\end{enumerate}
\end{lemma}

\begin{proof}
Let $\Lambda=\max_j|\lambda_j|$, $d=\max\{\deg P_j:|\lambda_j|=\Lambda\}$ and
$\mathcal J=\{j:|\lambda_j|=\Lambda,\ \deg P_j=d\}$. Write $\lambda_j=\Lambda e^{i\theta_j}$ for
$j\in\mathcal J$ (the $\theta_j$ are distinct modulo $2\pi$) and let $c_j\ne0$ be the
leading coefficient of $P_j$. Then, for $n\ge\max(n_0,1)$,
\[
  u_n=n^d\Lambda^n\,(v_n+\varepsilon_n),\qquad v_n=\sum_{j\in\mathcal J}c_je^{in\theta_j},\qquad
  \varepsilon_n\to0 .
\]
The sequences $n\mapsto n^e\lambda^n$ are linearly independent on every tail, so the
representation of $(u_n)$ is unique; as $u_n$ is real, the data $\{(\lambda_j,P_j)\}$ are
invariant under complex conjugation. Hence $\mathcal J$ is closed under conjugation, with
conjugate coefficients, and $v_n$ is real. By orthogonality of distinct characters
$n\mapsto e^{in\theta}$,
\[
  \lim_{N\to\infty}\frac1N\sum_{n<N}|v_n|^2=\sum_{j\in\mathcal J}|c_j|^2=:\sigma>0. \tag{$*$}
\]
In particular $|v_n|^2\ge\sigma/2$ for infinitely many $n$.

(a) If $\Lambda>1$, or $\Lambda=1$ and $d\ge1$, then $n^d\Lambda^n\to\infty$, and along the
infinitely many $n$ with $|v_n|\ge(\sigma/2)^{1/2}$ we get $|u_n|\to\infty$, contradicting
boundedness. Hence $\Lambda\le1$, and $d=0$ if $\Lambda=1$; this is (a).

(b) Now every $\theta_j\notin2\pi\Z$, hence $\frac1N\sum_{n<N}v_n\to0$. Suppose $u_n\ge0$ for
all $n\ge n_1$. Given $\varepsilon>0$ we then have $v_n\ge-\varepsilon$ for all large $n$, so
the upper mean of $v_n^-=\max(-v_n,0)$ is at most $\varepsilon$; since the mean of $v_n$ is
$0$, the upper mean of $v_n^+$ is also at most $\varepsilon$, and the upper mean of $|v_n|$
is at most $2\varepsilon$. As $|v_n|\le\sum|c_j|$, the upper mean of $|v_n|^2$ is at most
$2\varepsilon\sum|c_j|$. Letting $\varepsilon\to0$ contradicts $(*)$.
\end{proof}

%======================================================================
\section{Carry forcing}\label{sec:carry}
%======================================================================

\begin{thm}[Carry forcing for $0/1$-relations]\label{thm:carrygeneral}
Let $k\ge1$, $I\subseteq\{0,\dots,k-1\}$ with $|I|\ge2$, and let $\gamma>1$ be a root of
$F(x)=x^k-\sum_{i\in I}x^i$. Then:
\begin{enumerate}[label=(\roman*)]
\item $\gamma$ is the unique positive root of $F$, $1<\gamma<2$, and $\gamma$ is an
irrational algebraic integer;
\item for every $t>0$ and $n\ge0$ the carry
$\delta_n=s_{n+k}-\sum_{i\in I}s_{n+i}$ equals $\fl{\sum_{i\in I}f_{n+i}}\in\{0,\dots,|I|-1\}$;
\item for every $t>0$, $\delta_n\ge1$ for infinitely many $n$.
\end{enumerate}
\end{thm}

\begin{proof}
(i) Let $b=\min I$ and $F_0(x)=x^{-b}F(x)=x^{k-b}-\sum_{i\in I}x^{i-b}$; then $F_0(0)=-1$
and $F_0$ is monic with integer coefficients. The coefficient sequence of $F_0$ has exactly
one sign change, so by Descartes' rule of signs $F_0$ (hence $F$) has exactly one positive
root, counted with multiplicity. As $F_0(1)=1-|I|<0$ and
$F_0(2)=2^{k-b}-\sum_{i\in I}2^{i-b}\ge2^{k-b}-(2^{k-b}-1)=1$, this root is $\gamma\in(1,2)$.
Being a root of a monic integer polynomial, $\gamma$ is an algebraic integer; a rational
algebraic integer is an integer, and $1<\gamma<2$, so $\gamma\notin\Q$.

(ii) Multiplying \eqref{eq:rel} by $t\gamma^n$ gives
$s_{n+k}+f_{n+k}=\sum_{i\in I}(s_{n+i}+f_{n+i})$, so
$\delta_n=\sum_{i\in I}f_{n+i}-f_{n+k}$ is an integer in $(-1,|I|)$; since
$0\le f_{n+k}<1$ it equals $\fl{\sum_{i\in I}f_{n+i}}$.

(iii) Suppose $\delta_n=0$ for all $n\ge N$. Then $f_{n+k}=\sum_{i\in I}f_{n+i}$ for
$n\ge N$. Putting $g_m=f_m$ for $m\ge M:=N+b$ and substituting $n=m-b$,
\[
  g_{m+(k-b)}=\sum_{i\in I}g_{m+(i-b)}\qquad(m\ge M),
\]
a linear recurrence (involving only indices $\ge M$) with characteristic polynomial $F_0$,
whose roots are non-zero. Hence $g_m=\sum_{\lambda}P_\lambda(m)\lambda^m$ $(m\ge M)$, the sum
running over the distinct roots $\lambda$ of $F_0$, with polynomials $P_\lambda$ of degree less
than the multiplicity of $\lambda$. Discard the vanishing terms. Since $0\le g_m<1$,
Lemma~\ref{lem:exppoly}(a) shows that no remaining $\lambda$ has $|\lambda|>1$; in particular
the term at $\lambda=\gamma$ vanishes. The remaining roots are not positive reals, because
$\gamma$ is the only positive root of $F_0$. If some term remains,
Lemma~\ref{lem:exppoly}(b) yields $g_m<0$ for infinitely many $m$, which is absurd. Hence
$g_m=0$, i.e.\ $t\gamma^m\in\Z$, for all $m\ge M$. Then
$\gamma=t\gamma^{M+1}/t\gamma^M\in\Q$, contradicting (i).
\end{proof}

\begin{proof}[Proof of Theorem~\ref{thm:A}]
Apply Theorem~\ref{thm:carrygeneral} with $I=\{a,b\}$: the carries lie in $\{0,1\}$ and are
non-zero infinitely often.
\end{proof}

\begin{remark}[The golden ratio]\label{rem:phi}
For $\gamma=\varphi$ the argument is completely elementary: a carry-free tail gives
$f_{n+2}=f_{n+1}+f_n$, hence $f_n=A\varphi^n+B\psi^n$ with $\psi=-1/\varphi$; boundedness gives
$A=0$, and $B\psi^n\ge0$ for all large $n$ forces $B=0$, so $t\varphi^n\in\Z$ for all large $n$,
which is impossible. This is the argument of Kitamura \cite{Kit26}, and it is the only case of
Theorem~\ref{thm:A} needed for Theorem~\ref{thm:C} with $\gamma=\varphi$.
\end{remark}

\begin{example}
For $\gamma=\varphi$ and $t=1$ one has $\varphi^n=L_n-\psi^n$ with Lucas numbers $L_n$ and
$\psi=-1/\varphi$, so $\fl{\varphi^n}=L_n$ for odd $n$ and $\fl{\varphi^n}=L_n-1$ for even $n$.
Hence $\delta_n=1$ for every odd $n$ and $\delta_n=0$ for every even $n$. Numerically, for
the $28$ trinomial bases and seven values of $t$ tested in Section~\ref{sec:num} (check C2), the
proportion of $n<400$ with $\delta_n=1$ lies between $0.46$ and $0.54$, as one would expect if
$f_{n+a}$ and $f_{n+b}$ behaved like independent uniform variables.
\end{example}

\begin{remark}[Pisot trinomial bases]\label{rem:dub}
Let $\gamma$ be a Pisot number with $\gamma^k=\gamma^a+\gamma^b$, and put $F=x^k-x^a-x^b$. A
carry-free tail would force $\fr{t\gamma^n}\to0$: the bounded solution of the recurrence consists
of components with $|\lambda|\le1$; those with $|\lambda|=1$ vanish by the argument of
Lemma~\ref{lem:exppoly}(b) applied with $\Lambda=1$ (note $F(1)\ne0$), and the others tend to
$0$. On the other hand, by a theorem of Dubickas \cite[Theorem~4]{Dub06}, for an algebraic number
$\alpha>1$ there is $\xi>0$ such that $\fr{\xi\alpha^n}$ converges only if $\alpha$ is a strong Pisot
number (one whose largest conjugate other than $\alpha$ itself is positive), or $\alpha=2$, or
$\alpha$ is a Pisot number whose minimal polynomial $P$ satisfies $P(1)\le-2$. None of these
applies to $\gamma$: it has no positive conjugate other than itself, and $P(1)$ divides
$F(1)=-1$. Hence, for Pisot trinomial bases, Theorem~\ref{thm:A} also follows from
\cite{Dub06}. Our argument covers non-Pisot bases as well, e.g.\ $\sqrt\varphi$ and the roots of
$x^6-x^5-1$ and $x^5-x^3-1$ (Table~\ref{tab:roots}), and relations with $|I|\ge3$.
\end{remark}

\begin{remark}[Why the single positive root matters; Salem bases]\label{rem:salem}
The proof of (iii) uses only two facts: $\gamma\notin\Q$, and the relation polynomial has no
positive root other than $\gamma$. A Salem number $\gamma$ satisfies no relation
\eqref{eq:rel} at all: its minimal polynomial divides every integer polynomial vanishing at
$\gamma$ and also vanishes at $1/\gamma\in(0,1)$, whereas $F$ has the single positive root
$\gamma$. For a Salem base there are sequences $f_m=c\gamma^{-m}$, $c>0$, which are bounded
non-negative solutions of the recurrence given by the minimal polynomial, and Geneson's
counterexample \cite[Theorem~9]{Gen26} (all terms of $S_t(\gamma)$ even for suitable large
$t$) shows that completeness can indeed fail there. In particular, trinomial roots are never
Salem numbers.
\end{remark}

\begin{lemma}[The trinomial bases]\label{lem:trinroots}
The positive root of $x^k-x^a-x^b$ $(k>a>b\ge0)$ coincides with that of $x^m-x^j-1$, where
$m=k-b>j=a-b\ge1$. Let $\gamma_{m,j}$ denote the latter.
\begin{enumerate}[label=(\alph*)]
\item $\gamma_{2,1}=\varphi$ and $\gamma_{3,2}=1.46557\ldots$ exceed $\sqrt2$; for every other
pair $(m,j)$ we have $\gamma_{m,j}\le\gamma_{4,3}=1.38027\ldots<\sqrt2$, and
$\gamma_{m,j}\le\gamma_{3,1}=1.32471\ldots$ unless $(m,j)=(4,3)$.
\item $\gamma_{j+1,j}$ decreases strictly to $1$ as $j\to\infty$; and for every
$\varepsilon>0$ only finitely many pairs $(m,j)$ have $\gamma_{m,j}\ge1+\varepsilon$.
\end{enumerate}
\end{lemma}

\begin{proof}
The first assertion holds because $x^k-x^a-x^b=x^b(x^m-x^j-1)$. Write $\gamma=\gamma_{m,j}$,
so $\gamma^j(\gamma^{m-j}-1)=1$.
(a) If $m-j\ge2$ then $1\ge\gamma(\gamma^2-1)$, i.e.\ $\gamma^3-\gamma-1\le0$, so
$\gamma\le\gamma_{3,1}$. If $m-j=1$ then $\gamma^j(\gamma-1)=1$; for $j=1,2,3$ this gives
$\varphi$, $1.46557\ldots$, $1.38027\ldots$, and for $j\ge4$ we get $\gamma^4(\gamma-1)\le1$,
so $\gamma\le\gamma_{5,4}=\gamma_{3,1}$ (indeed $x^5-x^4-1=(x^2-x+1)(x^3-x-1)$). Finally,
$x^2-x-1$, $x^3-x^2-1$ and $x^4-x^3-1$ take the values $1-\sqrt2<0$, $2\sqrt2-3<0$ and
$3-2\sqrt2>0$ at $\sqrt2$, so $\gamma_{2,1},\gamma_{3,2}>\sqrt2>\gamma_{4,3}$.
(b) If $\gamma_{j+2,j+1}\ge\gamma_{j+1,j}=:\gamma'$, then
$\gamma_{j+2,j+1}^{\,j+1}(\gamma_{j+2,j+1}-1)\ge\gamma'^{\,j+1}(\gamma'-1)>\gamma'^{\,j}(\gamma'-1)=1$,
a contradiction. If $\gamma_{m,j}\ge1+\varepsilon$, then
$1=\gamma^j(\gamma^{m-j}-1)\ge(1+\varepsilon)^j\bigl((1+\varepsilon)^{m-j}-1\bigr)$, which
fails when $j$ or $m-j$ is large; so only finitely many $(m,j)$ qualify.
\end{proof}

\begin{table}[ht]
\centering
\small
\begin{tabular}{llll}
\toprule
$\gamma$ & trinomial(s) & Pisot & covered by \\
\midrule
$1.6180339887$ & $x^2-x-1$ & yes & Thm.~\ref{thm:C} (unions) \\
$1.4655712319$ & $x^3-x^2-1$ & yes & Thm.~\ref{thm:C} (unions) \\
$1.3802775691$ & $x^4-x^3-1$ & yes & Thms.~\ref{thm:B}, \ref{thm:C} \\
$1.3247179572$ & $x^3-x-1$, $x^5-x^4-1$ & yes & Thms.~\ref{thm:B}, \ref{thm:C} \\
$1.2851990332$ & $x^6-x^5-1$ & no & Thms.~\ref{thm:B}, \ref{thm:C} \\
$1.2720196495$ & $x^4-x^2-1$ \ ($\gamma=\sqrt\varphi$) & no & Thms.~\ref{thm:B}, \ref{thm:C} \\
$1.2554228711$ & $x^7-x^6-1$ & no & Thms.~\ref{thm:B}, \ref{thm:C} \\
$1.2365057034$ & $x^5-x^3-1$ & no & Thms.~\ref{thm:B}, \ref{thm:C} \\
$1.2207440846$ & $x^4-x-1$ & no & Thms.~\ref{thm:B}, \ref{thm:C} \\
$1.1938591113$ & $x^5-x^2-1$ & no & Thms.~\ref{thm:B}, \ref{thm:C} \\
$1.1673039783$ & $x^5-x-1$ & no & Thms.~\ref{thm:B}, \ref{thm:C} \\
$1.1347241384$ & $x^6-x-1$ & no & Thms.~\ref{thm:B}, \ref{thm:C} \\
\bottomrule
\end{tabular}
\caption{Some trinomial bases $\gamma_{m,j}$ (positive roots of $x^m-x^j-1$). By
Lemma~\ref{lem:trinroots}, all trinomial roots other than the first two are at most
$1.3803<\sqrt2$, and they accumulate only at $1$. Roots below the plastic number are never
Pisot numbers \cite{Sie44}.}
\label{tab:roots}
\end{table}

%======================================================================
\section{Unit switches and their extraction}\label{sec:extract}
%======================================================================

Fix a trinomial relation $\gamma^k=\gamma^a+\gamma^b$ with $k>a>b\ge0$ and $t>0$. For $n\ge0$
let $T_n=\{n+b,\,n+a,\,n+k\}$. If $\delta_n=1$ we call $T_n$ a \emph{unit switch}: the
disjoint index sets $\{n+b,n+a\}$ and $\{n+k\}$ carry sums $s_{n+a}+s_{n+b}$ and
$s_{n+a}+s_{n+b}+1$. For a set $E\subset\N$ we write
\[
  D_E(m)=s_m-\sum_{i<m,\ i\notin E}s_i\qquad(m\notin E),
\]
so that $D_\emptyset=D$; note that $D_E(m)$ increases when $E$ grows.

\begin{lemma}[Switch extraction]\label{lem:extract}
Assume moreover $\gamma<\sqrt2$. Let $\rho=\frac{2-\gamma}{\gamma-1}$ and
\[
  \kappa_1=\frac{\rho\gamma-1}{2\gamma},\qquad
  \kappa_2=\frac12\Bigl(\rho-\frac{\gamma+1}{\gamma^2}\Bigr)=\frac{1+\gamma^2-\gamma^3}{2\gamma^2(\gamma-1)},\qquad
  \kappa_3=\frac{\rho\gamma-2}{2\gamma}=\frac{2-\gamma^2}{2\gamma(\gamma-1)} .
\]
Then $\kappa_1,\kappa_2,\kappa_3>0$. Put $\kappa=\min(\kappa_1,\kappa_2,\kappa_3)$, choose an
integer $G\ge k+1$ such that $\eta:=2\gamma^{k-G}/(1-\gamma^{-G})$ satisfies
\begin{equation}\label{eq:eta}
  \eta\le\min\Bigl\{\tfrac12\gamma^{b}(\rho\gamma-1),\ \tfrac12\gamma^{a+1}\Bigl(\rho-\tfrac{\gamma+1}{\gamma^2}\Bigr),\ \tfrac12\gamma^{k}(\rho\gamma-2)\Bigr\},
\end{equation}
and choose $N_0$ such that $t\kappa\gamma^m>\frac t{\gamma-1}+m$ for all $m\ge N_0$. Let $R\ge1$,
$N_0\le n_1<n_2<\dots<n_R$ with $n_{j+1}-n_j\ge G$, and let $E=T_{n_1}\cup\dots\cup T_{n_R}$.
Then for every $m\notin E$
\[
  D_E(m)\ \le\ \max\bigl(c_0(t,\gamma),0\bigr).
\]
Here $G$ depends only on $(k,a,b)$, and $N_0$ only on $(k,a,b,t)$; neither depends on $R$ or on
the $n_j$.
\end{lemma}

\begin{proof}
\emph{Positivity of the constants.} $\rho\gamma-1=\frac{1+\gamma-\gamma^2}{\gamma-1}>0$ since
$\gamma<\varphi$; $1+\gamma^2-\gamma^3>0$ since $\gamma<1.4655\ldots$ (the real root of
$x^3-x^2-1$); $\rho\gamma-2=\frac{2-\gamma^2}{\gamma-1}>0$ since $\gamma<\sqrt2$. The identity
$\rho-\frac{\gamma+1}{\gamma^2}=\frac{\gamma^2(2-\gamma)-(\gamma^2-1)}{\gamma^2(\gamma-1)}$ gives
the stated form of $\kappa_2$. A suitable $G$ exists because $\eta\to0$ as $G\to\infty$.

\emph{Set-up.} Since $G\ge k+1>k-b$, the triples are pairwise disjoint and ordered: every
element of $T_{n_l}$ is smaller than every element of $T_{n_{l+1}}$. By
Lemma~\ref{lem:growth} and $s_i\le t\gamma^i$,
\begin{equation}\label{eq:DE}
  D_E(m)=D(m)+\sum_{i\in E,\,i<m}s_i\le -t\rho\gamma^m+\frac t{\gamma-1}+m+t\sum_{i\in E,\,i<m}\gamma^i .
\end{equation}
If $m<n_1+b$ then no element of $E$ is smaller than $m$ and $D_E(m)=D(m)\le c_0$. Let now
$m>n_1+b$, $m\notin E$, and let $j$ be the largest index with $n_j+b<m$; then $m<n_{j+1}+b$ if
$j<R$. The triples $T_{n_l}$, $l<j$, lie entirely below $m$ and, by
$\gamma^a+\gamma^b=\gamma^k$, contribute
\[
  \sum_{l<j}\bigl(\gamma^{n_l+b}+\gamma^{n_l+a}+\gamma^{n_l+k}\bigr)=\sum_{l<j}2\gamma^{n_l+k}
  \le2\gamma^{n_j+k}\sum_{r\ge1}\gamma^{-rG}=\eta\gamma^{n_j};
\]
the triples $T_{n_l}$, $l>j$, lie entirely above $m$. Let
$W=\sum_{i\in T_{n_j},\,i<m}\gamma^i$. We claim
\begin{equation}\label{eq:claim}
  \rho\gamma^m-W-\eta\gamma^{n_j}\ \ge\ \kappa\gamma^m .
\end{equation}
\emph{Case 1: $n_j+b<m<n_j+a$.} Then $W=\gamma^{n_j+b}\le\gamma^{m-1}$ and
$\eta\gamma^{n_j}\le\eta\gamma^{-b}\gamma^{m-1}$, so the left side of \eqref{eq:claim} is at least
$\gamma^{m-1}(\rho\gamma-1-\eta\gamma^{-b})\ge\gamma^{m-1}\cdot\frac12(\rho\gamma-1)=\kappa_1\gamma^m$
by \eqref{eq:eta}.

\emph{Case 2: $n_j+a<m<n_j+k$.} Then $W=\gamma^{n_j+b}+\gamma^{n_j+a}=\gamma^{n_j+k}$ and
$\gamma^{n_j}\le\gamma^{m-a-1}$. Moreover $\gamma^{k-a}=1+\gamma^{b-a}\le1+\gamma^{-1}$, so
$\gamma^{k-a-1}\le(\gamma+1)/\gamma^2$. Hence the left side of \eqref{eq:claim} is at least
$\gamma^m\bigl(\rho-\frac{\gamma+1}{\gamma^2}-\eta\gamma^{-a-1}\bigr)\ge\kappa_2\gamma^m$.

\emph{Case 3: $m>n_j+k$.} Then $W=2\gamma^{n_j+k}\le2\gamma^{m-1}$ and
$\eta\gamma^{n_j}\le\eta\gamma^{-k}\gamma^{m-1}$, so the left side is at least
$\gamma^{m-1}(\rho\gamma-2-\eta\gamma^{-k})\ge\kappa_3\gamma^m$.

Since $m\notin E$ these cases are exhaustive (some may be empty). By \eqref{eq:DE} and
\eqref{eq:claim}, $D_E(m)\le-t\kappa\gamma^m+\frac t{\gamma-1}+m<0$, because
$m>n_1+b\ge N_0$.
\end{proof}

\begin{proposition}[Sharpness of $\sqrt2$ for extraction]\label{prop:sharp}
Let $\gamma\in(\sqrt2,2)$ satisfy $\gamma^k=\gamma^a+\gamma^b$ and let $t>0$. For every $n\ge0$
and every $E\supseteq T_n$ with $n+k+1\notin E$,
\[
  D_E(n+k+1)\ \ge\ t\,\frac{\gamma^2-2}{\gamma-1}\,\gamma^{n+k}-4 .
\]
Consequently, for every $\gamma\in(\sqrt2,2)$ the conclusion of Lemma~\ref{lem:extract} fails
for every choice of $G$ and $N_0$: no bound $D_E(m)\le C(t,\gamma)$ can hold uniformly over
the admissible sets $E$, even over single triples $E=T_n$ with $\delta_n=1$. Among trinomial
roots this concerns exactly $\varphi$ and $1.46557\ldots$ (Lemma~\ref{lem:trinroots}).
\end{proposition}

\begin{proof}
By monotonicity in $E$ it suffices to take $E=T_n$. With $m=n+k+1$: $s_m\ge t\gamma^m-1$,
$\sum_{i<m}s_i\le t(\gamma^m-1)/(\gamma-1)$, and
$\sum_{i\in E}s_i\ge t(\gamma^{n+b}+\gamma^{n+a}+\gamma^{n+k})-3=2t\gamma^{n+k}-3$. Thus
\[
  D_E(m)\ge t\gamma^m-1-\frac{t(\gamma^m-1)}{\gamma-1}+2t\gamma^{n+k}-3
  =t\gamma^{n+k}(2-\rho\gamma)+\frac t{\gamma-1}-4,
\]
and $2-\rho\gamma=\frac{\gamma^2-2}{\gamma-1}>0$. The right-hand side tends to infinity with
$n$, and by Theorem~\ref{thm:A} there are switches $T_n$ with $n$ arbitrarily large.
\end{proof}

%======================================================================
\section{Proof of Theorem~B}\label{sec:proofB}
%======================================================================

\begin{proof}[Proof of Theorem~\ref{thm:B}]
Let $\gamma\in(1,\sqrt2)$ satisfy $\gamma^k=\gamma^a+\gamma^b$ and let $t>0$. Put
$c'=\max(c_0(t,\gamma),0)$ (Lemma~\ref{lem:growth}) and $R=c'+1$. Let $G,N_0$ be as in
Lemma~\ref{lem:extract}. By Theorem~\ref{thm:A} there are infinitely many $n$ with
$\delta_n=1$; choose $n_1<\dots<n_R$ among them with $n_1\ge N_0$ and $n_{j+1}-n_j\ge G$, and
let $E=\bigcup_jT_{n_j}$.

\emph{A run of length $R+1$.} For each $j$ select either $\{n_j+b,n_j+a\}$ or $\{n_j+k\}$.
These selections are disjoint, and the total sum equals $u+r$, where
$u=\sum_j(s_{n_j+a}+s_{n_j+b})$ and $r$ is the number of $j$ for which $\{n_j+k\}$ was
selected. Hence the subset sums of $\mathcal P=(s_i)_{i\in E}$ contain the $L:=R+1$
consecutive integers $u,\dots,u+R$.

\emph{Propagation.} Let $(q_m)$ enumerate $(s_i)_{i\notin E}$ in increasing order of $i$.
Then \eqref{eq:propcond} for $(q_m)$ reads $D_E(i)\le L$ for $i\notin E$, which holds by
Lemma~\ref{lem:extract} since $D_E(i)\le c'<L$. Also $\sum_mq_m=\infty$. By
Lemma~\ref{lem:prop}, every integer $N\ge u$ is a sum of distinct terms of
$(\fl{t\gamma^n})_{n\ge0}$.

\emph{Strong completeness.} Let $\Omega\subset\N$ be finite and delete the terms with indices in
$\Omega$. The growth bound becomes
$D_\Omega(m)\le-t\rho\gamma^m+\frac t{\gamma-1}+m+\sum_{i\in\Omega}s_i$, so
$c_0^\Omega:=\max_{m\notin\Omega}D_\Omega(m)<\infty$. Repeating the proof of
Lemma~\ref{lem:extract} with $E\cup\Omega$ in place of $E$, $c_0^\Omega$ in place of $c_0$, and
$N_0$ enlarged so that $N_0>\max\Omega$ and
$t\kappa\gamma^m>\frac t{\gamma-1}+m+\sum_{i\in\Omega}s_i$ for $m\ge N_0$, we obtain
$D_{E\cup\Omega}(m)\le\max(c_0^\Omega,0)$ for $m\notin E\cup\Omega$. Choosing
$R=\max(c_0^\Omega,0)+1$ switches
beyond $N_0$, the argument above applies verbatim.

\emph{The set of values.} Since $s_{n+1}-s_n\ge t\gamma^n(\gamma-1)-1>0$ for large $n$, only
finitely many indices $i$ satisfy $s_i=s_{i'}$ for some $i'<i$. Deleting these finitely many
indices leaves a complete family whose values are distinct and exhaust the set
$\{s_n:n\ge0\}$; hence this set is complete.
\end{proof}

\begin{remark}[Splitting]\label{rem:split}
Let $\gamma$ be the root of $x^m-x^j-1$ and $d=\gcd(m,j)\ge2$. Then $\beta=\gamma^d$ is the root of
the trinomial $y^{m/d}-y^{j/d}-1$, and $S_t(\gamma)$ contains the two families
$(\fl{t\beta^n})_{n\ge0}$ and $(\fl{(t\gamma)\beta^n})_{n\ge0}$ (indices $\equiv0$ and $\equiv1$
modulo $d$) on disjoint index sets. By Theorem~\ref{thm:C} their union is (strongly) complete,
hence so is $S_t(\gamma)$; for $\gamma=\sqrt\varphi$ this is the argument of \cite{Kit26}. When
$\gcd(m,j)=1$, as for $x^3-x-1$, $x^4-x^3-1$ and all $x^m-x^{m-1}-1$, no such splitting into
orbits of a trinomial base is available from the trinomial itself, and the proof above uses
Lemma~\ref{lem:extract} instead.
\end{remark}

\begin{remark}
The last step matters for formal versions of Problem \#349 that speak about the set of values
$\{\fl{t\gamma^n}:n\ge0\}$ (as the Formal Conjectures file for \#349 does \cite{FC}) rather than
about the indexed sequence. The threshold $u$ produced by the proof is enormous: it involves
$R=c_0+1$ switches, where $c_0=\fl t$ as soon as $t(2-\gamma)>2$ (by Bernoulli's inequality,
$D(m)\le t-t\rho(\gamma^m-1)+m<t-1$ for $m\ge1$),
switches, while numerically the sequences become complete much earlier (Section~\ref{sec:num}).
The proof is effective once the positions of unit carries are known; Theorem~\ref{thm:A}
itself gives no bound on the gaps between carries.
\end{remark}

%======================================================================
\section{Proof of Theorem~C and of the corollaries}\label{sec:proofC}
%======================================================================

\begin{proof}[Proof of Theorem~\ref{thm:C}]
Let $\gamma\in(1,2)$ be a trinomial base, $\alpha,\beta>0$, $A=(x_n)=(\fl{\alpha\gamma^n})_{n\ge0}$,
$B=(y_n)=(\fl{\beta\gamma^n})_{n\ge0}$. By Lemma~\ref{lem:growth},
$c_A:=\max(\max_m(x_m-\sum_{i<m}x_i),0)<\infty$. By Theorem~\ref{thm:A} applied to $t=\beta$,
the sequence $B$ has infinitely many unit switches; choose $R=c_A+1$ of them,
$T_{n_1},\dots,T_{n_R}$, with $n_{j+1}-n_j\ge k+1$ so that they are disjoint. As in the proof
of Theorem~\ref{thm:B}, the subset sums of these $3R$ terms of $B$ contain $L=R+1$
consecutive integers $u,\dots,u+R$. The sequence $(q_m)=A$ satisfies \eqref{eq:propcond}
with this $L$, because $x_m-\sum_{i<m}x_i\le c_A<L$. Since the terms of $A$ and $B$ carry
different indices, Lemma~\ref{lem:prop} shows that every integer $N\ge u$ is a sum of
distinct terms of $A\sqcup B$. Strong completeness: deleting finitely many terms of $A$
changes $c_A$ into a larger but finite constant, and the switches can be chosen beyond the
deleted terms of $B$.
\end{proof}

No term is removed from the propagating sequence $A$; this is why Theorem~\ref{thm:C} needs no
restriction $\gamma<\sqrt2$.

\subsection*{The formal statement of Problem \#354(ii)}
In \cite{FC} (file \nolinkurl{FormalConjectures/ErdosProblems/354.lean}; part (ii) has used
a general base $\gamma$ since a correction of September 1, 2026), one defines
$\mathtt{FloorMultiples}(a,\gamma)(n)=\fl{\gamma^na}\in\Z$ and
\[
  \mathtt{interleave}(\alpha,\beta,\gamma)(n)=
  \begin{cases}\fl{\gamma^{n/2}\alpha}, & n\ \text{even},\\
  \fl{\gamma^{(n-1)/2}\beta}, & n\ \text{odd},\end{cases}
\]
and the theorem \texttt{erdos\_354.parts.ii} (tagged \texttt{research open}) states
\[
  \exists\,\gamma\in(1,2)\ \forall\,\alpha>0\ \forall\,\beta>0:\ \ \alpha/\beta\notin\Q\
  \Longrightarrow\ \mathtt{IsAddCompleteNatSeq'}\bigl(\mathtt{interleave}(\alpha,\beta,\gamma)\bigr),
\]
where $\mathtt{IsAddCompleteNatSeq'}(A)$ means that every sufficiently large integer equals
$\sum_{i\in J}A(i)$ for some finite set $J\subset\N$ of indices.

\begin{proof}[Proof of Corollary~\ref{cor:D}]
The interleaved sequence is exactly the indexed union of $A=(\fl{\alpha\gamma^n})$ and
$B=(\fl{\beta\gamma^n})$, and $\mathtt{IsAddCompleteNatSeq'}$ is completeness of this indexed
family. Let $\gamma\in(1,2)$ be a trinomial root, e.g.\ $\gamma=\varphi$. By Theorem~\ref{thm:C}
the union is complete for all $\alpha,\beta>0$, in particular whenever $\alpha/\beta\notin\Q$.
(For $\gamma=\sqrt\varphi$ this is the theorem of \cite{Kit26}.) For $\gamma<\sqrt2$, already the set of values of $A$ is
complete by Theorem~\ref{thm:B}, so the union is complete also as a set of values.
\end{proof}

\begin{remark}[Multiset versus set]\label{rem:multiset}
Problem \#354 speaks of a \emph{multiset}, and \cite{FC} formalises completeness with sums over
sets of indices, so a value occurring in both sequences may be used twice. For $\gamma=\varphi$
this matters: with $\beta=25.5$ and $\alpha=\beta\varphi$ (so $\alpha/\beta=\varphi\notin\Q$) the
two sequences coincide up to a shift, the multiset union is complete by Theorem~\ref{thm:C},
but the set of values is that of $S_{25.5}(\varphi)$, which is not complete by \cite{vD26}
since $25.5\ge3/\varphi^2$. For the trinomial roots $\gamma<\sqrt2$ the distinction disappears.
\end{remark}

\begin{proof}[Proof of Corollary~\ref{cor:E}]
Apply Theorem~\ref{thm:C} with $\gamma=\varphi$ and $\alpha=\beta=t$. For the negative statement,
van Doorn \cite{vD26} shows that for $\varphi\le\gamma\le5^{1/3}$ the sequence $S_t(\gamma)$ is
complete if and only if $t<\min(3/\gamma^2,5/\gamma^3)$, and
$\min(3/\varphi^2,5/\varphi^3)=3/\varphi^2$.
\end{proof}

%======================================================================
\section{Discussion}\label{sec:disc}
%======================================================================

\subsection{Base \texorpdfstring{$2$}{2} (Problem \#354(i))}
Binary expansions are unique, so $\sum_{i\in I}2^i\ne\sum_{j\in J}2^j$ for distinct finite
sets $I\ne J$; in particular there is no relation \eqref{eq:rel} for $\gamma=2$, and the real
numbers $2^i\alpha$ $(i\ge0)$ satisfy no relation $\sum_{i\in I}2^i\alpha=\sum_{j\in J}2^j\alpha$
with $I\ne J$, so a single base-$2$ sequence has no carries of the kind used here. Between two
base-$2$ sequences, an exact relation
$\sum_{i}\varepsilon_i2^i\alpha=\sum_{j}\varepsilon'_j2^j\beta$ with
$\varepsilon_i,\varepsilon'_j\in\{0,\pm1\}$ forces $\alpha/\beta\in\Q$ unless both sides vanish.
Hence, when $\alpha/\beta\notin\Q$, any ``switch'' must come from an approximate relation, whose
carry is not controlled. This is the structural reason why part (i) resists the present
method, and why the irrationality of $\alpha/\beta$ plays such a different role in base $2$
(where $\beta=2^k\alpha$ can give incompleteness) and in base $\varphi$ (where
Theorem~\ref{thm:C} needs no hypothesis on $\alpha/\beta$).

\subsection{Salem bases and a dichotomy}
Remark~\ref{rem:salem} and Geneson's theorem \cite{Gen26} suggest a dichotomy among algebraic
bases: when a $0/1$-relation with a single positive root is available, carries are forced in
every orbit; for Salem numbers no such relation exists, the obstruction used here disappears, and
completeness can fail. It would be interesting to turn this into a characterisation.

\subsection{Open problems}
\begin{question}
Is $S_t(\gamma)$ complete for every $t>0$ when $\gamma=1.46557\ldots$ is the root of
$x^3-x^2-1$? (For $\gamma=\varphi$ the answer is negative for $t\ge3/\varphi^2$ \cite{vD26}.)
\end{question}

\begin{question}\label{q:units}
Let $\gamma$ satisfy a relation \eqref{eq:rel} with $|I|\ge3$ (for instance the tribonacci
constant, $\gamma^3=\gamma^2+\gamma+1$). Is $\delta_n(t)=1$ for infinitely many $n$, for every
$t>0$? An affirmative answer for such a $\gamma$ extends Theorem~\ref{thm:C} to $\gamma$ verbatim.
\end{question}

\begin{remark}
The roots of the polynomials $x^k-\sum_{i\in I}x^i$ are dense in $(1,2)$. This is classical
(simple Parry numbers are dense, see Parry \cite{Par60}); we include the short argument. Let
$\gamma_0\in(1,2)$ and let $1=\sum_{j\ge1}d_j\gamma_0^{-j}$ be its greedy expansion with digits
$d_j\in\{0,1\}$ (so $d_1=1$). The function $g_K(x)=1-\sum_{j\le K}d_jx^{-j}$ is increasing on
$[1,2]$ with $g_K'(x)\ge x^{-2}\ge\frac14$, and $0\le g_K(\gamma_0)\le\gamma_0^{-K}/(\gamma_0-1)$;
hence the root $\gamma_K$ of $x^Kg_K(x)=x^K-\sum_{j\le K}d_jx^{K-j}$ satisfies
$0\le\gamma_0-\gamma_K\le4\gamma_0^{-K}/(\gamma_0-1)\to0$ (and $|I|\ge2$ for large $K$, since
$g_K(\gamma_0)\to0$ while $1-\gamma_0^{-1}>0$). Thus a positive answer to
Question~\ref{q:units} for all relations with $|I|\ge3$ would give a dense set of bases for
Theorem~\ref{thm:C}.
\end{remark}

\begin{question}
Give an effective bound for the completeness threshold in Theorems~\ref{thm:B}
and~\ref{thm:C}, e.g.\ via a quantitative form of Lemma~\ref{lem:exppoly}(b) bounding the
gaps between carries.
\end{question}

\begin{question}
Characterise the set of $\gamma\in(1,2)$ for which the union in \eqref{eq:354ii} is complete
for all $\alpha,\beta>0$ (respectively, for all $\alpha,\beta$ with $\alpha/\beta\notin\Q$).
\end{question}

%======================================================================
\section{Numerical verification}\label{sec:num}
%======================================================================

The script in Appendix~\ref{app:script} computes floors $\fl{t\gamma^n}$ with \texttt{mpmath}
at two working precisions ($900$ and $1050$ significant digits, recomputing $\gamma$ and $t$)
and requires the integer outputs to agree. The finite computations illustrate the proofs;
the proofs do not depend on them. The script checks:
\begin{enumerate}[label=C\arabic*.,leftmargin=2.5em]
\item each of the $28$ trinomials $x^m-x^j-1$ with $1\le j<m\le8$ has a unique positive root,
which lies in $(1,2)$;
\item for these $28$ bases and seven values of $t$ (including $t=1,\sqrt2,\pi,1000.0001$) the
carries lie in $\{0,1\}$ and at least $50$ of the first $400$ are units;
\item the constants $\rho\gamma-1$, $\rho-(\gamma+1)/\gamma^2$, $\rho\gamma-2$ of
Lemma~\ref{lem:extract} are positive for the tested bases below $\sqrt2$, and $\rho\gamma-2<0$ for
$\varphi$ and $1.4655\ldots$;
\item finite sanity checks in the shape of the proof of Theorem~\ref{thm:B}: a choice of
$R=c'+1$ unit switches with gaps $\ge k+8$ (an ad hoc choice of $G$ and $N_0$), and a direct
check of $D_E(m)\le R+1$ for all $m\notin E$ up to $80$ beyond the last switch, for six bases
and $t\in\{1,2.5,17.3,123.456\}$; for instance, for the plastic number and $t=123.456$ one finds
$c_0=123$, $R=124$, and the last switch at index $1498$ (these are illustrations, not
certificates);
\item brute-force subset-sum computations up to $4\cdot10^5$ for four bases below $\sqrt2$
and $t\in\{1,2.5,7.77\}$: no non-representable integer lies in $[10^5,4\cdot10^5]$;
\item for $\gamma=\varphi$, the data of the proof of Theorem~\ref{thm:C} (the constant $c_A$ and
$R=c_A+1$ disjoint unit switches in $B$) for
$(\alpha,\beta)$ equal to $(3.7,3.7)$, $(1,\sqrt2)$, $(\pi,10)$ and $(50,50)$, and brute-force
tails for $\alpha=\beta=3.7$ and $\alpha=\beta=10$;
\item for $\varphi$ and $1.4655\ldots$ (and $t=2.5$), the deficit $D_{T_n}(n+k+1)$ is positive
and increases along $n=100$, $200$, $300$, $400$ (Proposition~\ref{prop:sharp}).
\end{enumerate}
All checks pass (total running time about $75$ seconds). The computations are consistent with
van Doorn's certified ranges \cite{vD26} and illustrate that the actual completeness thresholds
are far smaller than the ones produced by the proofs.

\appendix
\section{Verification script}\label{app:script}
The following Python script (also reproduced verbatim in this single \TeX\ file) performs the
checks C1--C7 of Section~\ref{sec:num}.

\begin{lstlisting}[language=Python]
#!/usr/bin/env python3
"""
Verification script for the paper
  "Carry Forcing for Zero-One Relations: the Erdos-Graham Completeness Prediction
   for All Trinomial Bases below sqrt 2".
Checks (exact integer arithmetic on floors computed at two working precisions):
  C0  prints the numerical values quoted in the paper (Table 1, x^10-x-1, 3/phi^2);
  C1  each trinomial x^m - x^j - 1 has a unique positive root gamma, and 1 < gamma < 2;
  C2  carries delta_n lie in {0,1}, and unit carries are frequent (Theorem A);
  C3  the three constants of the extraction lemma are positive iff gamma < sqrt 2;
  C4  finite sanity checks in the shape of the proof of Theorem B (not certificates);
  C5  brute-force completeness tails for small t (single sequences, gamma < sqrt 2);
  C6  the data of the proof of Theorem C for the golden ratio, including an
      instance of Erdos #354(ii) with alpha/beta irrational;
  C7  sharpness: for gamma > sqrt 2 removing a switch triple creates a growing deficit.
The finite computations illustrate the proofs; the proofs do not depend on them.
Requires: python3, mpmath.
"""
from mpmath import mp, mpf, findroot, floor, sqrt, pi, polyroots

mp.dps = 900
EXTRA = 150          # second, higher working precision used to confirm every floor


def troot(m, j):
    """unique positive root of x^m - x^j - 1 at the current precision"""
    x0 = mpf('1.6') if (m, j) == (2, 1) else mpf('1.3')
    return findroot(lambda x: x**m - x**j - 1, x0)


def floors(tmk, gmk, N):
    """floor(t*g^n) for n < N; t = tmk(), g = gmk() are recomputed at two precisions
    and the two lists of floors must coincide."""
    res = []
    for extra in (0, EXTRA):
        with mp.workdps(mp.dps + extra):
            t, g = tmk(), gmk()
            res.append([int(floor(t * g**n)) for n in range(N)])
    assert res[0] == res[1], "floor not determined at working precision"
    return res[0]


def const(v):
    return lambda: mpf(v)


BASES = [(m, j) for m in range(2, 9) for j in range(1, m)]   # all 28 trinomials, m <= 8
TS = [const(1), const('2.5'), (lambda: sqrt(2)), (lambda: +pi), const('17.3'),
      const('123.456'), const('1000.0001')]


def table_C0():
    """values quoted in the paper: trinomial roots (Table 1, x^10-x-1) and 3/phi^2"""
    for (m, j) in sorted(BASES + [(10, 1)], key=lambda p: -float(troot(*p))):
        print("    root of x^%d - x^%d - 1 = %s" % (m, j, mp.nstr(troot(m, j), 12)))
    phi = troot(2, 1)
    print("    3/phi^2 = %s,  5/phi^3 = %s" % (mp.nstr(3 / phi**2, 8), mp.nstr(5 / phi**3, 8)))
    print("C0 ok: table of trinomial roots printed")


def check_C1():
    for (m, j) in BASES:
        g = troot(m, j)
        assert abs(g**m - g**j - 1) < mpf(10) ** (-500)
        assert 1 < g < 2
        coeffs = [0] * (m + 1); coeffs[0] = 1; coeffs[m - j] -= 1; coeffs[m] -= 1
        rts = polyroots(coeffs, maxsteps=300, extraprec=300)
        pos = [r for r in rts if abs(r.imag) < mpf(10) ** (-30) and r.real > 0]
        assert len(pos) == 1 and abs(pos[0].real - g) < mpf(10) ** (-30)
    print("C1 ok: unique positive root in (1,2) for every tested trinomial")


def check_C2():
    fr = []
    for (m, j) in BASES:
        k, a, b = m, j, 0
        for tmk in TS:
            s = floors(tmk, lambda: troot(m, j), 400 + k)
            d = [s[n + k] - s[n + a] - s[n + b] for n in range(400)]
            assert set(d) <= {0, 1}
            assert sum(d) >= 50, (m, j, sum(d))
            fr.append(sum(d) / 400)
    print("    proportion of n < 400 with delta_n = 1: min %.3f, median %.3f, max %.3f"
          % (min(fr), sorted(fr)[len(fr) // 2], max(fr)))
    print("C2 ok: delta_n in {0,1}; unit carries frequent for all tested (base, t)")


def lemma_constants(g):
    rho = (2 - g) / (g - 1)
    return rho * g - 1, rho - (g + 1) / g**2, rho * g - 2


def check_C3():
    for (m, j) in BASES:
        g = troot(m, j)
        c1, c2, c3 = lemma_constants(g)
        if g < sqrt(2):
            assert c1 > 0 and c2 > 0 and c3 > 0
        else:
            assert c3 < 0
    print("C3 ok: extraction constants positive exactly for the tested bases below sqrt 2")


def sigma_constant(s):
    pref, c0 = 0, None
    for x in s:
        D = x - pref
        c0 = D if c0 is None else max(c0, D)
        pref += x
    return c0


def proof_shape_check_single(m, j, tmk, Nmax=2600):
    k, a, b = m, j, 0
    s = floors(tmk, lambda: troot(m, j), Nmax)
    c0 = sigma_constant(s)
    R = max(c0, 0) + 1                      # number of unit switches; interval length R+1
    G, N0 = k + 8, 2 * k + 5
    sw, last = [], -10**9
    for n in range(N0, Nmax - k - 1):
        if s[n + k] - s[n + a] - s[n + b] == 1 and n - last >= G:
            sw.append(n); last = n
            if len(sw) == R:
                break
    assert len(sw) == R, "not enough switches within the horizon"
    E = set()
    for n in sw:
        E |= {n + b, n + a, n + k}
    pref = 0
    for mm in range(sw[-1] + k + 80):
        if mm in E:
            continue
        D = s[mm] - pref                    # D_E(mm)
        assert D <= R + 1, ("check fails", m, j, mm, D, R)
        pref += s[mm]
    return c0, R, sw[-1]


def check_C4():
    for (m, j) in [(3, 1), (4, 3), (4, 2), (4, 1), (5, 3), (6, 5)]:
        for v in ['1', '2.5', '17.3', '123.456']:
            c0, R, last = proof_shape_check_single(m, j, const(v))
            if (m, j, v) == (3, 1, '123.456'):
                print("    plastic number, t = 123.456: c0 = %d, R = %d, last switch = %d" % (c0, R, last))
    print("C4 ok: sanity checks in the shape of the proof of Theorem B passed")


def tail_misses(seq, upto):
    S, mask = 1, (1 << (upto + 1)) - 1
    for x in seq:
        if 0 < x <= upto:
            S |= (S << x) & mask
    return [i for i in range(upto + 1) if not (S >> i) & 1]


def check_C5():
    upto = 400000
    for (m, j) in [(3, 1), (4, 3), (4, 2), (4, 1)]:
        for v in ['1', '2.5', '7.77']:
            s = floors(const(v), lambda: troot(m, j), 120)
            miss = tail_misses(s, upto)
            assert (miss[-1] if miss else -1) < upto // 4, (m, j, v)
    print("C5 ok: no non-representable integers in [upto/4, upto] for the tested small t")


def check_C6():
    gmk = lambda: troot(2, 1)
    pairs = [(const('3.7'), const('3.7')), (const(1), lambda: sqrt(2)),
             ((lambda: +pi), const('10')), (const('50'), const('50'))]
    for (amk, bmk) in pairs:
        A = floors(amk, gmk, 1500); B = floors(bmk, gmk, 1500)
        R = max(sigma_constant(A), 0) + 1
        sw, last = [], -10
        for n in range(len(B) - 3):
            if B[n + 2] - B[n + 1] - B[n] == 1 and n - last >= 3:
                sw.append(n); last = n
                if len(sw) == R:
                    break
        assert len(sw) == R
    for v in ['3.7', '10']:
        A = floors(const(v), gmk, 60)
        miss = tail_misses(A + A, 300000)
        assert (miss[-1] if miss else -1) < 75000
    print("C6 ok: golden-ratio unions: data of the proof of Theorem C found; brute-force tails consistent")


def check_C7():
    for (m, j) in [(2, 1), (3, 2)]:
        k, a, b = m, j, 0
        s = floors(const('2.5'), lambda: troot(m, j), 600)
        defs = []
        for n0 in [100, 200, 300, 400]:
            E = {n0 + b, n0 + a, n0 + k}
            mm = n0 + k + 1
            defs.append(s[mm] - sum(s[i] for i in range(mm) if i not in E))
        assert 0 < defs[0] < defs[1] < defs[2] < defs[3]
    print("C7 ok: for the bases above sqrt 2, a removed switch triple leaves a growing deficit")


if __name__ == "__main__":
    table_C0(); check_C1(); check_C2(); check_C3(); check_C4(); check_C5(); check_C6(); check_C7()
    print("ALL CHECKS PASSED")
\end{lstlisting}

\begin{thebibliography}{99}

\bibitem{BG07} J.~P.~Bell and S.~Gerhold, On the positivity set of a linear recurrence
sequence, \emph{Israel J. Math.} \textbf{157} (2007), 333--345.

\bibitem{Bro61} J.~L.~Brown, Jr., Note on complete sequences of integers, \emph{Amer. Math.
Monthly} \textbf{68} (1961), 557--560.

\bibitem{Dub06} A.~Dubickas, On the limit points of the fractional parts of powers of Pisot
numbers, \emph{Arch. Math. (Brno)} \textbf{42} (2006), 151--158.

\bibitem{EG80} P.~Erd\H{o}s and R.~L.~Graham, \emph{Old and New Problems and Results in
Combinatorial Number Theory}, Monographies de L'Enseignement Math\'ematique \textbf{28},
Universit\'e de Gen\`eve, 1980.

\bibitem{EP349} T.~F.~Bloom, Erd\H{o}s Problem \#349,
\url{https://www.erdosproblems.com/349}, accessed October 2026.

\bibitem{EP354} T.~F.~Bloom, Erd\H{o}s Problem \#354,
\url{https://www.erdosproblems.com/354}, accessed October 2026 (with references to
Graham, Hegyv\'ari, Jiang--Ma, Fang--He and van Doorn).

\bibitem{FC} Google DeepMind, \emph{Formal Conjectures},
\url{https://github.com/google-deepmind/formal-conjectures}, commit \texttt{5a2fb85}
(October 28, 2026); files \nolinkurl{FormalConjectures/ErdosProblems/349.lean} and
\nolinkurl{FormalConjectures/ErdosProblems/354.lean} (theorem \texttt{erdos\_354.parts.ii}).

\bibitem{Gen26} J.~Geneson, Deletion thresholds and exponential examples for complete
sequences, arXiv:2609.25107 (2026).

\bibitem{Gra64} R.~L.~Graham, Complete sequences of polynomial values, \emph{Duke Math. J.}
\textbf{31} (1964), 275--285.

\bibitem{He89} N.~Hegyv\'ari, Some remarks on a problem of Erd\H{o}s and Graham, \emph{Acta Math.
Hungar.} \textbf{53} (1989), 149--154 (not seen by us; quoted in \cite{Rev26}).

\bibitem{Kit26} K.~Kitamura, \emph{A Lean proof of Erd\H{o}s Problem 354(ii)}, GitHub repository
\url{https://github.com/KitaKen1/erdos-354-part-ii}, commit \texttt{5536b18} (September 5, 2026);
proposed to \cite{FC} as pull request \#5286.

\bibitem{Par60} W.~Parry, On the $\beta$-expansions of real numbers, \emph{Acta Math. Acad. Sci.
Hungar.} \textbf{11} (1960), 401--416.

\bibitem{Rev26} Automated review records of Formal Conjectures statements, GitHub repository
\url{https://github.com/tadamcz/fc-review-results}, record
\nolinkurl{data/aea251bb26/files/ErdosProblems/354.json} (September 23, 2026); see also
formal-conjectures issue \#6542 (September 24, 2026).

\bibitem{Sie44} C.~L.~Siegel, Algebraic integers whose conjugates lie in the unit circle,
\emph{Duke Math. J.} \textbf{11} (1944), 597--602.

\bibitem{vD26} W.~van Doorn, Completeness of exponentially increasing sequences,
arXiv:2602.23394 (2026).

\bibitem{YC26} Y.~Yu and K.~Chen, \emph{Erd\H{o}s Problem 354(i): Strong completeness of two
dyadic floor sequences}, manuscript and Lean formalization,
\url{https://github.com/Andrewyzzz/erdos354} (September 2026); proposed to \cite{FC} as pull
request \#6433.

\end{thebibliography}

\end{document}
